% ═══════════════════════════════════════════════════════════════════
%  PART 1: Introduction & KRLS Theory
% ═══════════════════════════════════════════════════════════════════

\section{Introduction}

\begin{tipbox}[Who Is This Guide For?]
This guide is written for \textbf{researchers, graduate students, and practitioners} in economics, finance, and environmental science who want to apply advanced nonparametric methods without requiring deep machine-learning expertise. We start from basic regression concepts and build up to the full QQKRLS estimator.
\end{tipbox}

\subsection{Motivation: Why Go Beyond OLS?}

Ordinary Least Squares (OLS) is the workhorse of empirical economics. Given a model $y = X\beta + \varepsilon$, OLS provides the best linear unbiased estimator (BLUE) under the Gauss--Markov assumptions. However, OLS imposes strong restrictions:

\begin{enumerate}[label=\textcolor{primaryblue}{\arabic*.}]
\item \textbf{Linearity}: The conditional expectation $\E[y \mid x]$ must be linear in $x$.
\item \textbf{Additivity}: Each predictor's effect is separable --- no automatic interactions.
\item \textbf{Homogeneity}: The effect $\beta_j$ is \emph{constant} across the entire distribution.
\end{enumerate}

In reality, economic relationships are often \textbf{nonlinear}, \textbf{interactive}, and \textbf{heterogeneous across the distribution}. For example, the effect of energy consumption on environmental quality may differ at low vs.\ high consumption levels.

\begin{warningbox}[The Problem with Misspecification]
If the true data-generating process is nonlinear but we fit a linear model, OLS estimates are \textbf{biased and inconsistent}. Diagnostic tests (Ramsey RESET, BDS) can detect this misspecification, but the solution requires a more flexible estimator.
\end{warningbox}

\subsection{The QQKRLS Framework}

The \texttt{qqkrls} library provides a unified solution through four methods:

\begin{center}
\begin{tabular}{>{\bfseries\color{deepblue}}l p{10cm}}
\toprule
\textbf{Method} & \textbf{Description} \\
\midrule
KRLS & Kernel Regularized Least Squares --- nonparametric regression in a reproducing kernel Hilbert space \\
QQ Regression & Quantile-on-Quantile analysis --- examines how effects vary across both distributions \\
WQR & Weighted Quantile Regression --- kernel-weighted quantile estimates \\
QQKRLS & The combined estimator --- nonparametric marginal effects across all quantile pairs \\
\bottomrule
\end{tabular}
\end{center}


% ═══════════════════════════════════════════════════════════════════
\section{Kernel Regularized Least Squares (KRLS)}
% ═══════════════════════════════════════════════════════════════════

\subsection{Intuition: From Linear to Nonparametric}

Consider the classical regression problem: given data $\{(x_i, y_i)\}_{i=1}^N$ where $x_i \in \R^D$ and $y_i \in \R$, we want to estimate the conditional expectation function $f(x) = \E[y \mid x]$.

\begin{itemize}[leftmargin=2em]
\item \textbf{OLS}: Restricts $f(x) = x^\top \beta$ --- a linear function.
\item \textbf{KRLS}: Allows $f$ to be \emph{any smooth function} in a reproducing kernel Hilbert space $\Hk$.
\end{itemize}

The key insight is the \textbf{kernel trick}: instead of explicitly specifying a functional form, KRLS defines similarity between observations using a \emph{kernel function} and learns a flexible, data-driven relationship.

\subsection{The Gaussian Kernel}

\begin{definitionbox}[Definition: Gaussian (RBF) Kernel]
The Gaussian kernel measures the similarity between two observations $x_i$ and $x_j$:
\begin{equation}
K(x_i, x_j) = \exp\!\left(-\frac{\|x_i - x_j\|^2}{\sigma}\right)
\label{eq:kernel}
\end{equation}
where $\sigma > 0$ is the \textbf{bandwidth} parameter controlling the smoothness of the function.
\end{definitionbox}

\textbf{Properties:}
\begin{itemize}[leftmargin=2em]
\item $K(x_i, x_i) = 1$ --- each observation is perfectly similar to itself.
\item $K(x_i, x_j) \to 0$ as $\|x_i - x_j\| \to \infty$ --- distant observations have negligible influence.
\item \textbf{Large $\sigma$}: Broader influence $\to$ smoother function (closer to linear).
\item \textbf{Small $\sigma$}: Narrower influence $\to$ more flexible (risk of overfitting).
\end{itemize}

\begin{tipbox}[Default Bandwidth]
Following Hainmueller \& Hazlett (2014), KRLS sets $\sigma = D$ (the number of predictors). This heuristic works well in practice and avoids expensive bandwidth search.
\end{tipbox}

\subsection{The KRLS Optimization Problem}

\begin{theorembox}[KRLS Objective Function]
KRLS solves the following Tikhonov-regularized optimization:
\begin{equation}
\hat{f} = \arg\min_{f \in \Hk} \underbrace{\sum_{i=1}^{N} (y_i - f(x_i))^2}_{\text{Data fidelity}} + \underbrace{\lambda \|f\|_{\Hk}^2}_{\text{Regularization}}
\label{eq:krls_obj}
\end{equation}
where $\lambda > 0$ controls the trade-off between fitting the data and keeping the function smooth.
\end{theorembox}

\textbf{Interpretation:}
\begin{itemize}[leftmargin=2em]
\item The first term minimizes the sum of squared residuals (like OLS).
\item The second term penalizes complexity in the RKHS norm, preventing overfitting.
\item $\lambda \to 0$: Interpolates the data (overfitting).
\item $\lambda \to \infty$: Converges to the sample mean (underfitting).
\end{itemize}

\subsection{The Representer Theorem and Closed-Form Solution}

\begin{theorembox}[Representer Theorem (Scholkopf \& Smola, 2002)]
The minimizer of Eq.~\eqref{eq:krls_obj} has the finite representation:
\begin{equation}
\hat{f}(x) = \sum_{i=1}^{N} c_i \cdot K(x, x_i)
\label{eq:representer}
\end{equation}
where $\bc = (c_1, \ldots, c_N)^\top$ is a vector of $N$ coefficients.
\end{theorembox}

This is remarkable: although the RKHS is infinite-dimensional, the optimal solution only depends on $N$ parameters. The predicted value at any point $x$ is a \textbf{weighted sum of kernel similarities} to all training observations.

\subsubsection{Eigendecomposition Solution}

Substituting Eq.~\eqref{eq:representer} into Eq.~\eqref{eq:krls_obj} and taking the first-order condition:
\begin{equation}
\bc^* = (\bK + \lambda \bI)^{-1} \by
\label{eq:coef_basic}
\end{equation}
where $\bK$ is the $N \times N$ kernel matrix with $K_{ij} = K(x_i, x_j)$.

For numerical stability, KRLS uses the eigendecomposition $\bK = \bV \bLambda \bV^\top$:
\begin{equation}
\boxed{\bc^* = \bV (\bLambda + \lambda \bI)^{-1} \bV^\top \by}
\label{eq:coef_eigen}
\end{equation}

\begin{examplebox}[Computational Advantage]
The eigendecomposition is computed \textbf{once} and reused for different values of $\lambda$, making leave-one-out cross-validation extremely efficient.
\end{examplebox}

\subsection{Regularization Parameter Selection via LOO-CV}

\begin{definitionbox}[Leave-One-Out Cross-Validation]
The LOO-CV error is:
\begin{equation}
\text{LOO}(\lambda) = \sum_{i=1}^{N} \left(\frac{y_i - \hat{f}_\lambda(x_i)}{1 - H_{ii}(\lambda)}\right)^2
\end{equation}
where $H_{ii}(\lambda)$ is the $i$-th diagonal of the hat matrix $\bH = \bK(\bK + \lambda \bI)^{-1}$.
\end{definitionbox}

KRLS selects $\lambda^*$ by minimizing this LOO error, which can be computed \textbf{without actually refitting} the model $N$ times (using the eigendecomposition).

\subsection{Pointwise Marginal Effects}

Unlike OLS (which gives one constant coefficient per variable), KRLS provides a \textbf{marginal effect for each observation}. This is computed analytically:

\begin{theorembox}[KRLS Pointwise Marginal Effect]
The marginal effect of predictor $d$ at observation $j$ is:
\begin{equation}
\frac{\partial \hat{f}(x_j)}{\partial x_j^{(d)}} = \frac{-2}{\sigma} \sum_{i=1}^{N} c_i \cdot K(x_j, x_i) \cdot \left(x_i^{(d)} - x_j^{(d)}\right)
\label{eq:deriv}
\end{equation}
\end{theorembox}

\textbf{Key insight}: The marginal effect varies by observation. If these pointwise effects show substantial variation (different at Q25 vs.\ Q75), the relationship is genuinely nonlinear, justifying KRLS over OLS.

\begin{examplebox}[Interpreting Marginal Effects]
\begin{itemize}[leftmargin=1.5em]
\item \textbf{Average Marginal Effect}: $\overline{\text{ME}}_d = \frac{1}{N}\sum_{j=1}^{N} \frac{\partial \hat{f}}{\partial x_j^{(d)}}$ --- comparable to OLS $\hat{\beta}_d$.
\item \textbf{Heterogeneity}: If Q25 $\neq$ Q75 of the effects, the relationship is nonlinear.
\item \textbf{Plots}: Scatter plots of marginal effects vs.\ predictor values reveal the functional form.
\end{itemize}
\end{examplebox}
